% revision_notice_v2.tex
% Front-matter revision notice, introduced in v2.0.
% Input from main.tex IMMEDIATELY BEFORE the abstract (author's instruction,
% handoff 2026-07-20-paper02-revision-notice-text; procedure per handoff
% 2026-07-18-papers-i-ii-zenodo-reversion-procedure).
%
% The notice paragraph below is the author's VERBATIM approved wording.
% Do not paraphrase, re-order, or "improve" it without author sign-off.
% The changes table follows on the next page.
%
% Structure deliberately mirrors Paper I's sections/revision_notice_v2.tex
% so the two revised papers read as a matched pair.

\medskip
\noindent\textbf{Revision notice, July 2026.}  This paper identified
the $(1, 1, -1)$ optical axis but did not follow its consequences
through the induced gauge-sector response.  Subsequent work shows
that the original single-domain three-diagonal hop set yields an
order-unity common-mode photon-speed anisotropy and is excluded as a
physical vacuum construction.  Results depending on that specific
gauge-sector realization should therefore be read conditionally.  See
\href{https://doi.org/10.5281/zenodo.21435952}%
{DOI: 10.5281/zenodo.21435952}~\cite{menendez2026opticalaxis}.

\newpage

\noindent\textbf{Changes from the previous version.}  The following
table records the effect of the result cited above on the claims of
versions v1.0--v1.01 of this paper.  Rows marked \emph{Withdrawn} are
withdrawn \emph{as constructed}: they concern the specific
three-diagonal single-domain hop set, not the $\mathcal{A}=1$
principle itself.  Rows inherited from Paper~I are reproduced here
because Paper~II builds on that construction.  Paper~IV gives the
full account.

\begin{center}
\footnotesize
\renewcommand{\arraystretch}{1.15}
\begin{tabular}{@{}p{0.30\textwidth} p{0.17\textwidth} p{0.45\textwidth}@{}}
\toprule
\textbf{\rule{0pt}{2.4ex}Original claim} &
\textbf{New status} &
\textbf{Reason\rule[-1.2ex]{0pt}{0pt}} \\
\midrule
Three-diagonal $\mathcal{A}=1$ vacuum is phenomenologically viable
  & Withdrawn
  & Gauge common-mode speed is $O(1)$ anisotropic.  Inherited from
    Paper~I; it is the substrate on which this paper's automorphism
    analysis is carried out \\
Matter sector safe (dimension-six-suppressed)
  & Withdrawn / retracted
  & Referee M1/M2 (verified in Paper~IV,
    \nolinkurl{exp_05_matter_sector_order}): matter also excluded as
    constructed.  A Paper~I/IV claim, listed for completeness \\
Optical axis is a suppressed lattice artifact
  & Revised
  & It enters the gauge action at dimension four.  This paper names
    the axis $\mathbf{V}_1 + \mathbf{V}_2 + \mathbf{V}_3 = (1,1,-1)$
    in \S\ref{sec:conclusion} but treats its consequence as a
    structural corollary rather than an exclusion risk \\
Bipartite-plaquette birefringence along $(1,1,-1)$ offered as a
testable consequence
  & Revised
  & The $Q$-tensor with eigenvalues $\{4, 4, 16\}$ is inherited from
    Paper~I unchanged; it is the same anisotropy now excluded as
    constructed.  The non-abelian corollary stands as algebra but is
    not a viable physical prediction of this vacuum \\
Substrate geometry of the derivation
  & Superseded
  & The three-diagonal hop set predicts an optical axis that is
    observable but is not observed, so the lattice geometry must
    change; the successor architecture is not birefringent.  Paper~IV
    gives the full account \\
$\mathcal{A}=1$ conservation principle
  & Unresolved, not falsified by this result
  & Failure traces to hop-set realization.  The Lie-algebra
    characterisation of Eq.~(137) --- containment $\supseteq$, the
    dim-71 centralizer, the factor-product projection --- is algebraic
    and independent of hop-set viability \\
\bottomrule
\end{tabular}
\end{center}

\newpage
